\section*{Bab \ref{ch:g11:trigo} --- Trigonometri: Lingkaran Satuan}

\begin{solution}{exo:g11:trigo:1}
Ke \omterm{def:g11:trigo:radian}{radian}: $30^\circ = \frac{\pi}{6}$, $45^\circ = \frac{\pi}{4}$,
$120^\circ = \frac{2\pi}{3}$, $270^\circ = \frac{3\pi}{2}$.

Ke derajat: $\frac{\pi}{5} = 36^\circ$, $\frac{5\pi}{6} = 150^\circ$,
$\frac{7\pi}{4} = 315^\circ$, $\frac{2\pi}{9} = 40^\circ$.
\end{solution}

\begin{solution}{exo:g11:trigo:2}
$\frac{2\pi}{3}$ dan $-\frac{\pi}{4}$ sudah berada di
$\intoc{-\pi}{\pi}$: berturut-turut di kuadran kedua dan kuadran keempat.

$\frac{17\pi}{6} - 2\pi = \frac{5\pi}{6}$: kuadran kedua, dekat dengan
sumbu $x$ negatif.

$-\frac{7\pi}{2} + 4\pi = \frac{\pi}{2}$: puncak lingkarannya, yaitu
$(0, 1)$.
\end{solution}

\begin{solution}{exo:g11:trigo:3}
$\cos\frac{3\pi}{4} = \cos\left(\pi - \frac\pi4\right)
= -\cos\frac\pi4 = -\frac{\sqrt2}{2}$.

$\sin\frac{5\pi}{6} = \sin\left(\pi - \frac\pi6\right)
= \sin\frac\pi6 = \frac12$.

$\cos\left(-\frac\pi3\right) = \cos\frac\pi3 = \frac12$.

$\sin\frac{4\pi}{3} = \sin\left(\pi + \frac\pi3\right)
= -\sin\frac\pi3 = -\frac{\sqrt3}{2}$.

$\cos\frac{11\pi}{6} = \cos\left(-\frac\pi6 + 2\pi\right)
= \cos\frac\pi6 = \frac{\sqrt3}{2}$.
\end{solution}

\begin{solution}{exo:g11:trigo:4}
$\sin^2 t = 1 - \cos^2 t = 1 - \frac{25}{169} = \frac{144}{169}$, sehingga
$\sin t = \pm\frac{12}{13}$. Pada $\intoo{-\frac\pi2}{0}$ (kuadran keempat)
\omterm{def:g10:coordgeom:system}{ordinatnya} negatif: jadi $\sin t = -\frac{12}{13}$.
\end{solution}

\begin{solution}{exo:g11:trigo:5}
Dengan memakai \cref{prop:g11:trigo:associated}:
\[
A = (-\cos x) + \cos x + \cos x - \cos x = 0 ,
\]
\[
B = \sin x + (-\sin x) + (-\sin x) = -\sin x .
\]
\end{solution}

\begin{solution}{exo:g11:trigo:6}
$\cos t = \frac{\sqrt3}{2}$: satu penyelesaiannya $\frac\pi6$, sehingga
$t = \pm\frac\pi6$ (pergeseran sebesar $2k\pi$ keluar dari \omterm{def:g10:numbers:interval}{intervalnya}).

$\sin t = \frac12$: $t = \frac\pi6$ atau $t = \pi - \frac\pi6 =
\frac{5\pi}{6}$.

$\cos t = -1$: satu-satunya titik yang berabsis $-1$ adalah $M(\pi)$, jadi $t = \pi$.
\end{solution}

\begin{solution}{exo:g11:trigo:7}
$\sin t = -\frac{\sqrt3}{2}$: dari $a = -\frac\pi3$, keluarganya adalah
$-\frac\pi3 + 2k\pi$ dan $\pi + \frac\pi3 + 2k\pi$; di
$\intco{0}{2\pi}$: $t = \frac{4\pi}{3}$ dan $t = \frac{5\pi}{3}$.

$\cos t = 0$: $t = \frac\pi2$ dan $t = \frac{3\pi}{2}$.

$2\sin t + 1 = 0$ berarti $\sin t = -\frac12$: dari $a = -\frac\pi6$, di
$\intco{0}{2\pi}$: $t = \frac{7\pi}{6}$ dan $t = \frac{11\pi}{6}$.
\end{solution}

\begin{solution}{exo:g11:trigo:8}
Menurut \cref{thm:g11:trigo:equations}, $\cos 2t = \cos t$ berarti
$2t = t + 2k\pi$ atau $2t = -t + 2k\pi$, yaitu $t = 2k\pi$ atau
$t = \frac{2k\pi}{3}$ ($k \in \Z$). Keluarga pertamanya termuat dalam
yang kedua (ambil $k$ kelipatan $3$), sehingga himpunan penyelesaiannya
\[
S = \left\{\frac{2k\pi}{3} : k \in \Z\right\}
\]
--- pada lingkarannya, yaitu ketiga titik $M(0)$,
$M\!\left(\frac{2\pi}{3}\right)$, $M\!\left(\frac{4\pi}{3}\right)$.
\end{solution}

\begin{solution}{exo:g11:trigo:9}
\omterm{def:g10:algebra:expand}{Menjabarkan} kedua kuadratnya lalu memakai $\cos^2 x + \sin^2 x = 1$:
\[
(\cos^2 x + 2\sin x\cos x + \sin^2 x)
+ (\cos^2 x - 2\sin x\cos x + \sin^2 x)
= 1 + 1 = 2 .
\]
\end{solution}

\begin{solution}{exo:g11:trigo:10}
Menggelinding tanpa tergelincir berarti panjang busurnya sama dengan jarak
yang ditempuh: $r\theta = 150$ cm memberi
$\theta = \frac{150}{30} = 5$ rad $= 5 \times \frac{180}{\pi} \approx
286.5^\circ$. Satu putaran penuh memajukan sepedanya sejauh kelilingnya,
$2\pi \times 30 = 60\pi \approx 188.5$ cm, sekitar $1.88$ m.
\end{solution}

\begin{solution}{exo:g11:trigo:11}
Pada lingkaran satuan, titik yang berabsis tepat $\frac12$ adalah
$M\!\left(\pm\frac\pi3\right)$. Titik yang berabsis \emph{paling besar}
$\frac12$ membentuk busur di sebelah kiri garis tegak $x = \frac12$,
yang dilalui ketika $t$ berjalan dari $\frac\pi3$ sampai $\pi$, atau dari
$-\pi$ sampai $-\frac\pi3$. Jadi, pada $\intoc{-\pi}{\pi}$:
\[
S = \intoc{-\pi}{-\frac{\pi}{3}} \cup \intcc{\frac{\pi}{3}}{\pi}.
\]
(Titik ujung $\pm\frac\pi3$ termasuk karena pertidaksamaannya takmurni.)
\end{solution}

\begin{solution}{pb:g11:trigo:1}
\textbf{1.} $30^\circ = \frac{\pi}{6}$; $135^\circ =
\frac{3\pi}{4}$; $300^\circ = \frac{5\pi}{3}$. Lalu
$\frac{3\pi}{4} = 135^\circ$; $\frac{5\pi}{6} = 150^\circ$.

\textbf{2.} Panjangnya $r\,t = 3 \times 2 = 6$ m. \omterm{def:g11:trigo:radian}{Radian} membuat panjang
busur menjadi hasil kali yang polos --- derajat akan menyeret
$\frac{\pi}{180}$ ke dalam setiap rumus.

\textbf{3.} \omterm{def:g11:trigo:radian}{Radian} yang ditempuh $= \frac{\text{jarak}}
{\text{jari-jari}} = \frac{1000}{0.30} \approx 3\,333$ rad;
dibagi $2\pi$: sekitar $530$ putaran penuh.

\textbf{4.} $\cos\frac{2\pi}{3} = -\frac12$;
$\sin\left(-\frac\pi4\right) = -\frac{\sqrt2}{2}$;
$\frac{19\pi}{6} - 2\pi = \frac{7\pi}{6}$, sehingga
$\cos\frac{19\pi}{6} = \cos\frac{7\pi}{6} =
-\frac{\sqrt3}{2}$.

\textbf{5.} $\cos^2 t = 1 - \frac{9}{25} = \frac{16}{25}$; pada
kuadran kedua kosinusnya negatif:
$\cos t = -\frac45$, dan
$\tan t = \frac{3/5}{-4/5} = -\frac34$.

\textbf{6.} Dalam $t$ menit kincirnya berputar
$\frac{t}{30} \times 2\pi = \frac{\pi t}{15}$. Di $t = 0$:
$h = 65 - 60 = 5$ m --- yaitu peron naik di bawah
(tinggi poros dikurangi jari-jari). Tanda minusnya menempatkanmu
\emph{di bawah} poros ketika sudut yang ditempuh $0$.

\textbf{7.} $h(7.5) = 65 - 60\cos\frac{\pi}{2} = 65$ m (setinggi
porosnya); $h(15) = 65 + 60 = 125$ m (puncaknya);
$h(20) = 65 - 60\cos\frac{4\pi}{3} = 65 + 30 = 95$ m.

\textbf{8.} $65 - 60\cos\frac{\pi t}{15} = 95$ memberi
$\cos\frac{\pi t}{15} = -\frac12$: penyelesaian pertamanya
$\frac{\pi t}{15} = \frac{2\pi}{3}$, yaitu $t = 10$ menit.

\textbf{9.} $\cos\frac{\pi t}{15} \leq -\frac12$ untuk
$\frac{\pi t}{15} \in \intcc{\frac{2\pi}{3}}{\frac{4\pi}{3}}$,
yaitu $t \in \intcc{10}{20}$: sepuluh menit perjalanan di atas
$95$ m, setangkup terhadap puncaknya di $t = 15$.

\textbf{10.} Menit pertama:
\[
h(1) - h(0) = 60\left(1 - \cos\frac{\pi}{15}\right)
\approx 1.3 \text{ m}.
\]
Antara menit $7$ dan $8$:
\[
60\left(\cos\frac{7\pi}{15} - \cos\frac{8\pi}{15}\right)
\approx 12.5 \text{ m}
\]
--- hampir sepuluh kali lebih banyak. Tingginya berubah paling cepat ketika
melewati ketinggian poros dan melambat di dekat dasar dan puncaknya:
sinusoidanya curam di tengah, mendatar di ujung ekstremnya.

\textbf{11.} Jari-jarinya $= \frac{40 - 4}{2} = 18$ m, porosnya di
$4 + 18 = 22$ m, kalanya $12$ menit:
$h(t) = 22 - 18\cos\left(\frac{\pi t}{6}\right)$.

\textbf{12.} \omterm{def:g10:functions:extrema}{Maksimumnya} $7 + 3 = 10$ m ketika
$\sin\frac{\pi t}{6} = 1$: pertama kali di $t = 3$ (pukul 03.00);
\omterm{def:g10:functions:extrema}{minimumnya} $4$ m pertama kali di $t = 9$; kalanya
$\frac{2\pi}{\pi/6} = 12$ jam.

\textbf{13.} $\sin\frac{\pi t}{6} \geq \frac12$ untuk
$\frac{\pi t}{6} \in \intcc{\frac\pi6}{\frac{5\pi}{6}}$:
$t \in \intcc{1}{5}$ --- jendela empat jam dari pukul 01.00 sampai
05.00, yang terbuka lagi satu kala kemudian, $t \in \intcc{13}{17}$
(pukul 13.00 sampai 17.00).

\textbf{14.} Cermin terhadap sumbu tegak mengirim titik bersudut
$t$ ke titik bersudut $\pi - t$, dan mempertahankan
\omterm{def:g10:coordgeom:system}{ordinatnya}: $\sin(\pi - t) = \sin t$. Setengah putaran terhadap
pusatnya mengirim $t$ ke $\pi + t$, dan meniadakan kedua \omterm{def:g10:coordgeom:system}{koordinatnya}:
$\cos(\pi + t) = -\cos t$. Lalu cermin terhadap diagonalnya
menukar \omterm{def:g10:coordgeom:system}{absis} dan \omterm{def:g10:coordgeom:system}{ordinat}:
$\cos\left(\frac\pi2 - t\right) = \sin t$ --- yaitu ``ko''
pada sudut yang saling berpenyiku.

\textbf{15.} $\sin(-t) = -\sin t$: ganjil (\omterm{def:g10:functions:graph}{grafiknya} setangkup
setengah putaran); $\cos(-t) = \cos t$: genap (setangkup cermin). Lalu
$\sin t + t^3$ adalah jumlah dua \omterm{def:g11:func:parity}{fungsi ganjil}: jadi ganjil.

\textbf{16.} Faktorkan: $(2\sin t + 1)(\sin t - 1) = 0$:
$\sin t = 1$ memberi $t = \frac\pi2$;
$\sin t = -\frac12$ memberi $t = \frac{7\pi}{6}$ dan
$t = \frac{11\pi}{6}$. Jadi tiga penyelesaian pada lingkarannya.

\textbf{17.} $1$ rad $= \frac{180}{\pi} \approx 57.3^\circ$.
Jadi $\sin(1) \approx 0.841$ sedangkan $\sin(1^\circ) \approx
0.0175$: hampir $50$ kali lipat. Moralnya: periksalah tanda modenya
sebelum memercayai tampilan trigonometri mana pun.

\textbf{18.} \omterm{def:g10:coordgeom:system}{Absis} sama dengan \omterm{def:g10:coordgeom:system}{ordinat} pada diagonalnya:
$t = \frac\pi4$ dan $t = \frac{5\pi}{4}$.

\textbf{19.} $\sin(0.1) = 0.09983\ldots$ terhadap $0.1$:
galat nisbinya sekitar $0.17\,\%$. Untuk sudut kecil, tali busurnya
memeluk busurnya --- dan limit $\frac{\sin t}{t} \to 1$ tahun
berikutnya persis merupakan pengamatan ini yang dijadikan
sebuah teorema.

\textbf{20.} \omterm{def:g11:trigo:radian}{Radian}: sudut $=$ busur pada lingkaran satuan, sehingga
sudut dihitung seperti panjang. \omterm{def:g11:trigo:cossin}{Kosinus} dan \omterm{def:g11:trigo:cossin}{sinus}: \omterm{def:g10:coordgeom:system}{koordinat}
jarum jamnya --- bayangan di tanah, tinggi di
dinding. Kesimetrian lingkarannya: seluruh katalog
identitas, yang terbaca dari cermin dan setengah putaran. \omterm{def:g10:algebra:equation}{Persamaan}
$\cos t = c$: jadwal jam itu, dengan lingkarannya memperlihatkan
semua penyelesaian sekaligus. Busananya: kincirnya mengenakan
$65 - 60\cos$, pasang surutnya mengenakan $7 + 3\sin$ --- jam yang sama,
matematika yang sama.

\end{solution}
